Considere un circuito ideal en serie con inductancia $L>0$, resistencia $R>0$, tensión aplicada $E(t)$ y corriente $I(t)$.

Los componentes son ideales, L y R son constantes positivas y se desprecia cualquier capacitancia. La orientación de corriente es la del generador. En (a) la tensión es constante y t=0 es el cierre; en (b) E es continua en el intervalo de tiempo considerado y puede variar.

\begin{enumerate}[label=(\alph*),ref=\alph*,leftmargin=*]
\item Una batería de $12\,\mathrm V$ se conecta a un circuito con $L=0.5\,\mathrm H$ y $R=10\,\Omega$. La corriente inicial es cero. Determine la corriente después de cerrar el circuito.
\item Para una tensión continua $E(t)$, compruebe que la solución general de $LI^{\prime}+RI=E(t)$ puede escribirse como \[I(t)=e^{-Rt/L}\left[C+\frac1L\int_0^t e^{Rs/L}E(s)\,ds\right].\] Especialice la fórmula a $E(t)=E_0$ constante e identifique el término transitorio y la corriente de régimen.
\end{enumerate}
